\documentclass[11pt]{article}
% Préambule commun aux documents ELECARM Crisis2 (généré via LaTeX / tectonic)
\usepackage[T1]{fontenc}
\usepackage[utf8]{inputenc}
\usepackage{lmodern}
\usepackage[french]{babel}
\usepackage{textcomp}
\usepackage[table]{xcolor}
\usepackage{tabularx}
\usepackage{array}
\usepackage[most]{tcolorbox}
\usepackage[a4paper,margin=2.2cm]{geometry}
\usepackage{parskip}
\usepackage{enumitem}
\usepackage{ragged2e}
\usepackage{helvet}

% ---- Palette ----
\definecolor{ink}{HTML}{111111}
\definecolor{brand}{HTML}{7F1D1D}
\definecolor{sub}{HTML}{555555}
\definecolor{ruleline}{HTML}{CCCCCC}
\definecolor{anomalybg}{HTML}{FBEAEA}
\definecolor{notebg}{HTML}{F8F4E8}
\definecolor{noteborder}{HTML}{B8860B}
\definecolor{quotebg}{HTML}{F2F2F2}
\definecolor{quoteborder}{HTML}{999999}
\definecolor{warnbg}{HTML}{FFF7ED}
\definecolor{warnborder}{HTML}{C2410C}
\definecolor{okbg}{HTML}{F0FDF4}
\definecolor{okborder}{HTML}{166534}
\definecolor{gouvblue}{HTML}{002654}
\definecolor{gouvred}{HTML}{ED2939}

\color{ink}
\renewcommand{\familydefault}{\rmdefault}
\setlength{\parindent}{0pt}
\linespread{1.12}

% ---- En-têtes (letterheads) ----
\newcommand{\lhrule}{\par\vspace{4pt}\noindent\textcolor{ink}{\rule{\linewidth}{1.4pt}}\par\vspace{14pt}}

\newcommand{\lhGouv}{%
  \noindent%
  \textcolor{gouvblue}{\rule{0.24\linewidth}{5pt}}%
  \textcolor{ruleline}{\rule{0.24\linewidth}{5pt}}%
  \textcolor{gouvred}{\rule{0.24\linewidth}{5pt}}\par\vspace{8pt}%
  {\sffamily\footnotesize\bfseries RÉPUBLIQUE FRANÇAISE}\\[1pt]%
  {\sffamily\scriptsize\itshape\textcolor{sub}{Liberté --- Égalité --- Fraternité}}\\[4pt]%
  {\sffamily\normalsize\bfseries MINISTÈRE DE L'INTÉRIEUR}\\[1pt]%
  {\sffamily\footnotesize\textcolor{sub}{Direction Générale de la Sécurité Intérieure}}%
  \lhrule
}

\newcommand{\lhElecarm}{%
  \noindent\colorbox{brand}{\makebox[14pt][c]{\rule{0pt}{12pt}}}\hspace{8pt}%
  \raisebox{2pt}{\parbox[b]{0.8\linewidth}{%
    {\sffamily\normalsize\bfseries ELECARM SAS}\\[1pt]%
    {\sffamily\footnotesize\textcolor{sub}{ZI Les Charmes --- 41000 Blois --- contact@elecarm.fr}}}}%
  \lhrule
}

\newcommand{\lhThales}{%
  \noindent{\sffamily\normalsize\bfseries THALES DEFENCE SYSTEMS}\\[1pt]%
  {\sffamily\footnotesize\textcolor{sub}{Division Systèmes de guidage --- Diffusion Restreinte}}%
  \lhrule
}

% ---- Titre + méta ----
\newcommand{\doctitle}[2]{%
  {\sffamily\LARGE\bfseries #1\par}\vspace{3pt}%
  {\sffamily\small\textcolor{sub}{#2}\par}%
  \vspace{6pt}\noindent\textcolor{ruleline}{\rule{\linewidth}{0.6pt}}\par\vspace{10pt}%
}

% ---- Section ----
\newcommand{\docsec}[1]{\par\vspace{8pt}{\sffamily\small\bfseries\textcolor{brand}{\MakeUppercase{#1}}}\par\vspace{4pt}}

% ---- Encadrés ----
\newtcolorbox{notebox}{colback=notebg,colframe=noteborder,boxrule=0pt,leftrule=3pt,arc=1pt,left=10pt,right=10pt,top=8pt,bottom=8pt,boxsep=1pt}
\newtcolorbox{quotebox}{colback=quotebg,colframe=quoteborder,boxrule=0pt,leftrule=3pt,arc=0pt,left=10pt,right=10pt,top=8pt,bottom=8pt,boxsep=1pt}
\newtcolorbox{warnbox}{colback=warnbg,colframe=warnborder,boxrule=0pt,leftrule=3pt,arc=1pt,left=10pt,right=10pt,top=8pt,bottom=8pt,boxsep=1pt}
\newtcolorbox{okbox}{colback=okbg,colframe=okborder,boxrule=0pt,leftrule=3pt,arc=1pt,left=10pt,right=10pt,top=8pt,bottom=8pt,boxsep=1pt}
\newtcolorbox{monobox}{colback=quotebg,colframe=quoteborder,boxrule=0.4pt,arc=1pt,left=10pt,right=10pt,top=8pt,bottom=8pt,boxsep=1pt}

% ---- Tables ----
\renewcommand{\arraystretch}{1.3}
\newcolumntype{Y}{>{\RaggedRight\arraybackslash}X}
\newcommand{\thd}[1]{{\sffamily\scriptsize\bfseries\textcolor{sub}{\MakeUppercase{#1}}}}
\newcommand{\anrow}{\rowcolor{anomalybg}}

% ---- Pied de page ----
\newcommand{\docfooter}{%
  \par\vspace{18pt}\noindent\textcolor{ruleline}{\rule{\linewidth}{0.4pt}}\par\vspace{3pt}%
  {\sffamily\scriptsize\textcolor{sub}{Document produit dans le cadre de l'exercice pédagogique ELECARM SAS --- Opération Silent Gateway --- usage interne exclusivement.}}%
}

\begin{document}
\lhGouv
\doctitle{Directive nationale de vigilance - Secteurs sensibles}{Circulaire MININT-DGSI-2026-0517 - 18 mai 2026 - Diffusion : sous-traitants industrie de la défense}

Dans le contexte géopolitique actuel et à l'approche des élections présidentielles, le niveau de vigilance des entreprises opérant dans des secteurs sensibles (défense, énergie, santé, infrastructures critiques) est relevé. Plusieurs incidents récents témoignent d'un intérêt accru d'acteurs hostiles pour ces organisations, avec un mode opératoire récurrent : l'exploitation de vulnérabilités humaines plutôt que de failles purement techniques.

\docsec{Obligations pour les sous-traitants de l'industrie de la défense}
Tout sous-traitant intervenant sur un programme classé Diffusion Restreinte doit transmettre, sous \textbf{15 jours} à compter de la réception de la présente circulaire, un rapport d'analyse de sécurité comportant :
\begin{enumerate}[leftmargin=1.4em,itemsep=2pt,topsep=4pt]
  \item Un diagnostic des vulnérabilités organisationnelles et techniques identifiées ;
  \item Un plan de préconisations correctives priorisées ;
  \item Un rappel de sensibilisation du personnel aux techniques de manipulation.
\end{enumerate}

\docsec{Rappel - les 4 leviers de manipulation humaine (VICE)}
\begin{tabularx}{\linewidth}{@{}p{3cm}Y@{}}
\thd{Levier} & \thd{Principe}\\[2pt]
\hline
Vanité & Flatter l'ego de la cible, lui proposer un statut ou une reconnaissance qu'elle n'a pas en interne.\\
Idéologie & Faire appel à des convictions personnelles (politiques, militantes) pour justifier une trahison.\\
Contrainte & Faire pression par la menace, le chantage ou l'exploitation d'une faute passée.\\
Espèces & Offrir un avantage financier direct pour compenser une difficulté matérielle ressentie par la cible.\\
\end{tabularx}

\vspace{10pt}
\begin{notebox}
\textbf{Rappel :} une cible n'est presque jamais recrutée par un seul levier isolé. C'est généralement la combinaison de plusieurs signaux (pression financière, valorisation personnelle, contexte professionnel) qui crée une fenêtre d'opportunité pour un acteur malveillant.
\end{notebox}
\docfooter
\end{document}
